% TODO checklist:
% - [x] Inspected src/security/tenants.py (Multi-tenancy implementation)
% - [x] Reviewed src/config.py (Tenant configuration settings)
% - [x] Reviewed Q&A.md (Multi-tenancy design patterns)
% - [x] Reviewed README.md (Tenant isolation features)
% - [x] Reviewed db/postgres_control/repositories (Tenant-scoped queries)

\section{Multi-Tenancy}
\label{sec:multi-tenancy}

The Cineca Agentic Platform is designed from the ground up as a multi-tenant system, enabling secure and isolated operation for multiple organizations, research groups, or departments within a single deployment. Multi-tenancy is implemented at multiple architectural layers to ensure complete data isolation while maximizing resource efficiency.

\subsection{Tenant Isolation Model}
\label{subsec:tenant-isolation-model}

The platform implements a \textit{shared-infrastructure, isolated-data} multi-tenancy model. This approach provides:

\begin{itemize}
    \item \textbf{Cost Efficiency}: A single deployment serves multiple tenants, reducing operational overhead.
    \item \textbf{Data Isolation}: Each tenant's data is logically separated and access is strictly controlled.
    \item \textbf{Configuration Independence}: Tenants can have customized settings, rate limits, and defaults.
    \item \textbf{Compliance Support}: Audit trails and access controls are maintained per-tenant.
\end{itemize}

\subsection{Tenant Identification}
\label{subsec:tenant-identification}

Tenant context is established at the beginning of each request through a prioritized extraction process:

\begin{enumerate}
    \item \textbf{HTTP Header}: The \texttt{X-Tenant-Id} header takes highest precedence.
    \item \textbf{Query Parameter}: The \texttt{tenant}, \texttt{tid}, or \texttt{tenant\_id} query parameters.
    \item \textbf{Token Claims}: The \texttt{tid} or \texttt{tenant\_id} claim from the JWT.
    \item \textbf{Default Tenant}: A configured fallback when no tenant is specified.
\end{enumerate}

\begin{lstlisting}[language=Python,caption={Tenant selection with priority-based extraction}]
def select_tenant(
    request: Request | None = None,
    user: Any | None = None,
    *,
    fallback_to_default: bool = True,
    set_context: bool = True,
) -> TenantContext:
    """Resolve tenant id using: header -> query -> user -> default."""
    tenant_id: str | None = None
    source = "none"
    
    if request is not None:
        # 1) Check header
        header = request.headers.get(_header_name())
        if header:
            tenant_id, source = header, "header"
        else:
            # 2) Check query params
            qp = request.query_params.get(_query_param())
            if qp:
                tenant_id, source = qp, "query"
    
    if not tenant_id and user is not None:
        # 3) Extract from user/token claims
        for key in ("tenant_id", "tid", "tenant", "org"):
            if hasattr(user, key) and getattr(user, key):
                tenant_id, source = str(getattr(user, key)), "user"
                break
    
    if not tenant_id and fallback_to_default:
        tenant_id, source = _default_tenant(), "default"
    
    return TenantContext(id=tenant_id, source=source, allowed=True)
\end{lstlisting}

\subsection{Tenant Identifier Validation}
\label{subsec:tenant-validation}

Tenant identifiers are validated against a strict format to prevent injection attacks and ensure consistent namespacing:

\begin{lstlisting}[language=Python,caption={Tenant identifier validation}]
TENANT_RE = re.compile(r"^[A-Za-z][A-Za-z0-9._-]{0,63}$")

def _validate(tenant_id: str) -> None:
    if not tenant_id:
        raise ValueError("tenant id must not be empty")
    if not TENANT_RE.fullmatch(tenant_id):
        raise ValueError(
            "tenant id must match ^[A-Za-z][A-Za-z0-9._-]{0,63}$"
        )
\end{lstlisting}

This validation ensures:
\begin{itemize}
    \item Identifiers start with a letter (avoiding numeric or special character prefixes).
    \item Only alphanumeric characters, dots, underscores, and hyphens are allowed.
    \item Maximum length of 64 characters to fit standard database column constraints.
\end{itemize}

\subsection{Tenant Allowlist Enforcement}
\label{subsec:tenant-allowlist}

When \texttt{TENANCY\_ENABLED} is set to \texttt{true}, the platform can enforce an allowlist of valid tenants:

\begin{lstlisting}[language=Python,caption={Tenant allowlist configuration}]
# Configuration options
TENANCY_ENABLED = True          # Enable tenant enforcement
TENANCY_DEFAULT = "default"     # Fallback tenant
TENANCY_ALLOWED = [             # Allowed tenant list
    "acme-corp",
    "research-lab",
    "cineca-hpc"
]
# Or use wildcard: TENANCY_ALLOWED = "*"
\end{lstlisting}

\begin{lstlisting}[language=Python,caption={Allowlist enforcement logic}]
def _is_allowed(tenant_id: str) -> bool:
    allowed = _allowed_set()
    if allowed is None or "*" in allowed:
        return True
    if len(allowed) == 0:
        return True  # No allowlist configured = allow any
    return tenant_id.lower() in allowed

def enforce_tenant(request: Request, user: Any | None = None) -> str:
    """Resolve and enforce tenant selection."""
    ctx = select_tenant(request, user=user)
    if _enabled() and ctx.id and not ctx.allowed:
        raise HTTPException(
            status.HTTP_403_FORBIDDEN,
            detail={"message": "Tenant not allowed", "tenant": ctx.id}
        )
    return ctx.id or ""
\end{lstlisting}

\subsection{Context Variable Propagation}
\label{subsec:tenant-context}

The current tenant is stored in a Python context variable, making it accessible throughout the request lifecycle without explicit parameter threading:

\begin{lstlisting}[language=Python,caption={Tenant context variable management}]
_current_tenant: contextvars.ContextVar[str | None] = \
    contextvars.ContextVar("current_tenant", default=None)

def set_current_tenant(tenant_id: str | None) -> None:
    """Set the current tenant in context variable."""
    _current_tenant.set(tenant_id)

def get_current_tenant() -> str | None:
    """Return the current tenant id from context (or None)."""
    return _current_tenant.get()
\end{lstlisting}

This design enables lower layers (repositories, services) to access tenant context without modifying function signatures.

\subsection{Database-Level Isolation}
\label{subsec:db-isolation}

All database queries are automatically scoped to the current tenant:

\paragraph{PostgreSQL Control Plane}
Repository queries include implicit tenant filtering:

\begin{lstlisting}[language=Python,caption={Tenant-scoped database query example}]
class TenantScopedRepository(BaseRepository[T]):
    async def list(self, tenant_id: str, **filters) -> list[T]:
        query = select(self.model).where(
            self.model.tenant_id == tenant_id
        )
        for key, value in filters.items():
            query = query.where(getattr(self.model, key) == value)
        return await self.session.execute(query)
\end{lstlisting}

\paragraph{Memgraph Graph Database}
Graph queries are filtered by tenant boundaries to prevent cross-tenant data access:

\begin{lstlisting}[language=cypher,caption={Tenant-scoped Cypher query}]
MATCH (n:User {tenant_id: $tenant_id})-[:WORKS_AT]->(i:Institution)
WHERE i.tenant_id = $tenant_id
RETURN n, i
\end{lstlisting}

\paragraph{Redis Data Plane}
Cache and queue keys are namespaced by tenant:

\begin{lstlisting}[language=Python,caption={Tenant-scoped key namespacing}]
def tenantize_key(key: str, tenant_id: str | None = None) -> str:
    """Prefix cache/DB key with tenant to avoid collisions.
    
    Example: tenantize_key("rate:count") -> "t:acme:rate:count"
    """
    t = tenant_id if tenant_id is not None else get_current_tenant()
    t = t or "global"
    return f"t:{t}:{key}"
\end{lstlisting}

\subsection{Per-Tenant Configuration}
\label{subsec:per-tenant-config}

Tenants can have individualized settings that override platform defaults:

\begin{itemize}
    \item \textbf{Rate Limits}: Custom request and token limits per tenant.
    \item \textbf{Default Models}: Preferred LLM providers and model instances.
    \item \textbf{Tool Access}: Enabled or restricted tool categories.
    \item \textbf{Safety Settings}: PII scrubbing levels and intent filter strictness.
    \item \textbf{Cost Budgets}: Monthly or per-request spending limits.
\end{itemize}

\begin{lstlisting}[language=Python,caption={Per-tenant configuration example}]
@dataclass
class TenantConfig:
    tenant_id: str
    rate_limit_rpm: int = 60
    rate_limit_tpm: int = 40000
    default_model: str = "gpt-4o-mini"
    max_concurrent_jobs: int = 5
    pii_mode: str = "redact"
    allowed_tools: list[str] = field(default_factory=list)
    budget_usd_monthly: float | None = None
\end{lstlisting}

\subsection{Tenant-Scoped Audit Logging}
\label{subsec:tenant-audit}

All audit events include tenant context for compliance and forensic analysis:

\begin{lstlisting}[language=Python,caption={Audit event with tenant context}]
audit_policy_decision(
    policy="tenancy",
    subject=principal,
    action="select",
    resource=f"tenant:{tenant_id}",
    allowed=allowed,
    reason="tenant not in allowlist" if not allowed else None,
    tenant_id=tenant_id,
    attributes={"enabled": _enabled(), "source": source}
)
\end{lstlisting}

\subsection{FastAPI Integration}
\label{subsec:tenant-fastapi}

The \texttt{require\_tenant} dependency ensures tenant context is established for protected endpoints:

\begin{lstlisting}[language=Python,caption={Tenant requirement in FastAPI routes}]
def require_tenant():
    """Ensure tenant is selected and (if enabled) allowed."""
    async def _dep(
        request: Request, 
        user=Depends(get_current_user)
    ):
        return enforce_tenant(request, user=user)
    return _dep

# Usage in routers
@router.get("/items", dependencies=[Depends(require_tenant())])
async def list_items(tenant: str = Depends(get_current_tenant)):
    key = tenantize_key("items:list", tenant)
    # ... tenant-scoped operation
\end{lstlisting}

\subsection{Cross-Tenant Operations}
\label{subsec:cross-tenant}

Administrative operations that span tenants require explicit \texttt{admin:all} scope and are logged with elevated scrutiny:

\begin{itemize}
    \item Tenant creation and deletion
    \item Cross-tenant reporting and analytics
    \item System-wide configuration changes
    \item Bulk operations affecting multiple tenants
\end{itemize}

These operations are protected by additional authorization checks and generate detailed audit trails for compliance review.

